\begin{tcolorbox}[title=Box 2: Therapeutic strategies targeting CAFs in breast cancer]

\textbf{1. CAF depletion:}
\begin{itemize}
  \item FAP-targeted antibodies (sibrotuzumab, bispecific FAP $\times$ PD-1)
  \item FAP-targeted CAR-T cells
  \item FAP-targeted radioligand therapy ($^{68}$Ga-FAPI, $^{177}$Lu-FAPI)
  \item Anti-CAF antibody--drug conjugates (anti-PDPN ADCs)
\end{itemize}

\textbf{2. CAF reprogramming:}
\begin{itemize}
  \item All-trans retinoic acid (ATTRA)
  \item Vitamin D receptor agonists (calcipotriol, paricalcitol)
  \item BET inhibitors (JQ1, I-BET762)
  \item TGF-$\beta$ pathway inhibitors
\end{itemize}

\textbf{3. Signalling pathway blockade:}
\begin{itemize}
  \item CXCR4/CXCL12 axis (plerixafor)
  \item IL-6/STAT3 axis (tocilizumab, JAK inhibitors)
  \item TGF-$\beta$ signalling (bintrafusp alfa, galunisertib, ATE-Trap)
  \item FGF/FGFR signalling (erdafitinib)
\end{itemize}

\textbf{4. Metabolic targeting:}
\begin{itemize}
  \item Glycolysis inhibition (2-deoxyglucose)
  \item Glutaminase inhibition (CB-839)
  \item Fatty acid synthase inhibition
\end{itemize}

\textbf{5. Combination strategies:}
\begin{itemize}
  \item CAF depletion + immunotherapy
  \item CAF reprogramming + chemotherapy
  \item Anti-TGF-$\beta$ + anti-PD-1
  \item Metabolic inhibitors + immune checkpoint blockade
\end{itemize}

\textbf{Key principle:} Effective CAF-targeted therapy must address the heterogeneity and redundancy of CAF functions. Single-agent approaches are likely to be insufficient; combination strategies targeting multiple CAF mechanisms simultaneously are more promising.

\end{tcolorbox}
